% Módulo 08 — Referências (R762+; DOIs verificados + preprints declarados)
\begin{thebibliography}{9}
\bibitem{moura2021} MOURA, L. de; ULLRICH, S. The Lean 4 theorem prover and programming language. \textit{Lecture Notes in Computer Science}, p. 625--635, 2021. DOI: \url{https://doi.org/10.1007/978-3-030-79876-5_37}. Verificada no Crossref.
\bibitem{mathlib2020} THE MATHLIB COMMUNITY. The lean mathematical library. In: \textit{Proceedings of CPP 2020}, p. 367--381. DOI: \url{https://doi.org/10.1145/3372885.3373824}. Verificada no Crossref.
\bibitem{vaswani2017} VASWANI, A. et al. Attention is all you need. In: \textit{NeurIPS 30}, Long Beach, 2017. p. 5998--6008. (Sem DOI; obra de conferência.)
\bibitem{xuan2025} XUAN, H.; YANG, B.; LI, X. Exploring the impact of temperature scaling in softmax for classification and adversarial robustness. \textit{Preprint} arXiv:2502.20604, 2025. Disponível em: \url{https://arxiv.org/abs/2502.20604}. Acesso em: 10 out. 2026. (Sem DOI verificado; citado como preprint.)
\bibitem{komisarenko2024} KOMISARENKO, V.; KULL, M. Improving calibration by relating focal loss, temperature scaling, and properness. \textit{Preprint} arXiv:2408.11598, 2024. Disponível em: \url{https://arxiv.org/abs/2408.11598}. Acesso em: 10 out. 2026. (Sem DOI verificado.)
\bibitem{xie2024} XIE, J. et al. Calibrating language models with adaptive temperature scaling. \textit{Preprint} arXiv:2409.19817, 2024. Disponível em: \url{https://arxiv.org/abs/2409.19817}. Acesso em: 10 out. 2026. (Sem DOI verificado.)
\bibitem{bagnara2019} BAGNARA, R. et al. Correct approximation of IEEE 754 floating-point arithmetic for program verification. \textit{Preprint} arXiv:1903.06119, 2019. Disponível em: \url{https://arxiv.org/abs/1903.06119}. Acesso em: 10 out. 2026. (Sem DOI verificado.)
\bibitem{murray2026} MURRAY, T. Lipschitz-based robustness certification under floating-point execution. \textit{Preprint} arXiv:2603.13334, 2026. Disponível em: \url{https://arxiv.org/abs/2603.13334}. Acesso em: 10 out. 2026. (Sem DOI verificado.)
\bibitem{agentbench2023} LIU, X. et al. AgentBench: evaluating LLMs as agents. \textit{Preprint} arXiv:2308.03688, 2023. Disponível em: \url{https://arxiv.org/abs/2308.03688}. Acesso em: 10 out. 2026. (Sem DOI verificado.)
\bibitem{perspectivegap2026} SUN, Y. et al. PerspectiveGap: a benchmark for multi-agent orchestration prompting. \textit{Preprint} arXiv:2606.08878, 2026. Disponível em: \url{https://arxiv.org/abs/2606.08878}. Acesso em: 10 out. 2026. (Sem DOI verificado.)
\bibitem{orchrm2026} TSANG, K. Y. et al. Reward modeling for multi-agent orchestration. \textit{Preprint} arXiv:2606.13598, 2026. Disponível em: \url{https://arxiv.org/abs/2606.13598}. Acesso em: 10 out. 2026. (Sem DOI verificado.)
\bibitem{shannon1948} SHANNON, C. E. A mathematical theory of communication. \textit{Bell System Technical Journal}, v. 27, n. 3, p. 379--423, 1948. DOI: \url{https://doi.org/10.1002/j.1538-7305.1948.tb01338.x}.
\bibitem{kullback1951} KULLBACK, S.; LEIBLER, R. A. On information and sufficiency. \textit{The Annals of Mathematical Statistics}, v. 22, n. 1, p. 79--86, 1951. DOI: \url{https://doi.org/10.1214/aoms/1177729694}.
\bibitem{jaynes1957} JAYNES, E. T. Information theory and statistical mechanics. \textit{Physical Review}, v. 106, n. 4, p. 620--630, 1957. DOI: \url{https://doi.org/10.1103/PhysRev.106.620}.
\bibitem{bridle1990} BRIDLE, J. S. Probabilistic interpretation of feedforward classification network outputs, with relationships to statistical pattern recognition. \textit{Neurocomputing}, p. 227--236, 1990. DOI: \url{https://doi.org/10.1007/978-3-642-76153-9_28}.
\bibitem{gneiting2007} GNEITING, T.; RAFTERY, A. E. Strictly proper scoring rules, prediction, and estimation. \textit{Journal of the American Statistical Association}, v. 102, n. 477, p. 359--378, 2007. DOI: \url{https://doi.org/10.1198/016214506000001437}.
\bibitem{luong2015} LUONG, T.; PHAM, H.; MANNING, C. D. Effective approaches to attention-based neural machine translation. In: \textit{EMNLP 2015}, p. 1412--1421. DOI: \url{https://doi.org/10.18653/v1/D15-1166}.
\bibitem{devlin2019} DEVLIN, J. et al. BERT: pre-training of deep bidirectional Transformers for language understanding. In: \textit{NAACL 2019}, p. 4171--4186. DOI: \url{https://doi.org/10.18653/v1/N19-1423}.
\bibitem{dai2019} DAI, Z. et al. Transformer-XL: attentive language models beyond a fixed-length context. In: \textit{ACL 2019}, p. 2978--2988. DOI: \url{https://doi.org/10.18653/v1/P19-1285}.
\bibitem{wang2018} WANG, A. et al. GLUE: a multi-task benchmark and analysis platform for natural language understanding. In: \textit{BlackboxNLP/EMNLP 2018}, p. 353--355. DOI: \url{https://doi.org/10.18653/v1/W18-5446}.
\bibitem{rajpurkar2016} RAJPURKAR, P. et al. SQuAD: 100,000+ questions for machine comprehension of text. In: \textit{EMNLP 2016}, p. 2383--2392. DOI: \url{https://doi.org/10.18653/v1/D16-1264}.
\bibitem{deng2009} DENG, J. et al. ImageNet: a large-scale hierarchical image database. In: \textit{CVPR 2009}, p. 248--255. DOI: \url{https://doi.org/10.1109/CVPR.2009.5206848}.
\bibitem{wolf2020} WOLF, T. et al. Transformers: state-of-the-art natural language processing. In: \textit{EMNLP 2020: System Demonstrations}, p. 38--45. DOI: \url{https://doi.org/10.18653/v1/2020.emnlp-demos.6}.
\bibitem{hoare1969} HOARE, C. A. R. An axiomatic basis for computer programming. \textit{Communications of the ACM}, v. 12, n. 10, p. 576--580, 1969. DOI: \url{https://doi.org/10.1145/363235.363259}.
\bibitem{demoura2008} DE MOURA, L.; BJØRNER, N. Z3: an efficient SMT solver. In: \textit{LNCS}, p. 337--340, 2008. DOI: \url{https://doi.org/10.1007/978-3-540-78800-3_24}.
\bibitem{meurer2017} MEURER, A. et al. SymPy: symbolic computing in Python. \textit{PeerJ Computer Science}, v. 3, e103, 2017. DOI: \url{https://doi.org/10.7717/peerj-cs.103}.
\bibitem{fagan1976} FAGAN, M. E. Design and code inspections to reduce errors in program development. \textit{IBM Systems Journal}, v. 15, n. 3, p. 182--211, 1976. DOI: \url{https://doi.org/10.1147/sj.153.0182}.
\bibitem{mccabe1976} MCCABE, T. J. A complexity measure. \textit{IEEE Transactions on Software Engineering}, v. SE-2, n. 4, p. 308--320, 1976. DOI: \url{https://doi.org/10.1109/TSE.1976.233837}.
\bibitem{chidamber1994} CHIDAMBER, S. R.; KEMERER, C. F. A metrics suite for object oriented design. \textit{IEEE Transactions on Software Engineering}, v. 20, n. 6, p. 476--493, 1994. DOI: \url{https://doi.org/10.1109/32.295895}.
\bibitem{parnas1986} PARNAS, D. L.; CLEMENTS, P. C. A rational design process: how and why to fake it. \textit{IEEE Transactions on Software Engineering}, v. SE-12, n. 2, p. 251--257, 1986. DOI: \url{https://doi.org/10.1109/TSE.1986.6312940}.
\bibitem{smith1980} SMITH, R. G. The contract net protocol: high-level communication and control in a distributed problem solver. \textit{IEEE Transactions on Computers}, v. C-29, n. 12, p. 1104--1113, 1980. DOI: \url{https://doi.org/10.1109/TC.1980.1675516}.
\bibitem{erman1980} ERMAN, L. D. et al. The Hearsay-II speech-understanding system: integrating knowledge to resolve uncertainty. \textit{ACM Computing Surveys}, v. 12, n. 2, p. 213--253, 1980. DOI: \url{https://doi.org/10.1145/356810.356816}.
\bibitem{goldberg1991} GOLDBERG, D. What every computer scientist should know about floating-point arithmetic. \textit{ACM Computing Surveys}, v. 23, n. 1, p. 5--48, 1991. DOI: \url{https://doi.org/10.1145/103162.103163}.
\bibitem{muller2018} MULLER, J.-M. et al. \textit{Handbook of Floating-Point Arithmetic}. 2. ed. Cham: Springer, 2018. DOI: \url{https://doi.org/10.1007/978-3-319-76526-6}.
\bibitem{harris2020} HARRIS, C. R. et al. Array programming with NumPy. \textit{Nature}, v. 585, p. 357--362, 2020. DOI: \url{https://doi.org/10.1038/s41586-020-2649-2}.
\bibitem{virtanen2020} VIRTANEN, P. et al. SciPy 1.0: fundamental algorithms for scientific computing in Python. \textit{Nature Methods}, v. 17, p. 261--272, 2020. DOI: \url{https://doi.org/10.1038/s41592-019-0686-2}.
\bibitem{basili1986} BASILI, V. R.; SELBY, R. W.; HUTCHENS, D. H. Experimentation in software engineering. \textit{IEEE Transactions on Software Engineering}, v. SE-12, n. 7, p. 733--743, 1986. DOI: \url{https://doi.org/10.1109/TSE.1986.6312975}.
\bibitem{wohlin2012} WOHLIN, C. et al. \textit{Experimentation in Software Engineering}. Heidelberg: Springer, 2012. DOI: \url{https://doi.org/10.1007/978-3-642-29044-2}.
\bibitem{runeson2009} RUNESON, P.; HÖST, M. Guidelines for conducting and reporting case study research in software engineering. \textit{Empirical Software Engineering}, v. 14, n. 2, p. 131--164, 2009. DOI: \url{https://doi.org/10.1007/s10664-008-9102-8}.
\bibitem{page2021} PAGE, M. J. et al. The PRISMA 2020 statement: an updated guideline for reporting systematic reviews. \textit{BMJ}, n71, 2021. DOI: \url{https://doi.org/10.1136/bmj.n71}.
\bibitem{wasserstein2016} WASSERSTEIN, R. L.; LAZAR, N. A. The ASA statement on p-values: context, process, and purpose. \textit{The American Statistician}, v. 70, n. 2, p. 129--133, 2016. DOI: \url{https://doi.org/10.1080/00031305.2016.1154108}.
\bibitem{cumming2014} CUMMING, G. The new statistics: why and how. \textit{Psychological Science}, v. 25, n. 1, p. 7--29, 2014. DOI: \url{https://doi.org/10.1177/0956797613504966}.
\bibitem{efron1979} EFRON, B. Bootstrap methods: another look at the jackknife. \textit{The Annals of Statistics}, v. 7, n. 1, 1979. DOI: \url{https://doi.org/10.1214/aos/1176344552}.
\bibitem{osc2015} OPEN SCIENCE COLLABORATION. Estimating the reproducibility of psychological science. \textit{Science}, v. 349, n. 6251, 2015. DOI: \url{https://doi.org/10.1126/science.aac4716}.
\bibitem{ioannidis2005} IOANNIDIS, J. P. A. Why most published research findings are false. \textit{PLoS Medicine}, v. 2, n. 8, e124, 2005. DOI: \url{https://doi.org/10.1371/journal.pmed.0020124}.
\bibitem{collberg2016} COLLBERG, C.; PROEBSTING, T. A. Repeatability in computer systems research. \textit{Communications of the ACM}, v. 59, n. 3, p. 62--69, 2016. DOI: \url{https://doi.org/10.1145/2812803}.
\bibitem{gebru2021} GEBRU, T. et al. Datasheets for datasets. \textit{Communications of the ACM}, v. 64, n. 12, p. 86--92, 2021. DOI: \url{https://doi.org/10.1145/3458723}.
\bibitem{mitchell2019} MITCHELL, M. et al. Model cards for model reporting. In: \textit{FAT* 2019}, p. 220--229. DOI: \url{https://doi.org/10.1145/3287560.3287596}.
\bibitem{wilkinson2016} WILKINSON, M. D. et al. The FAIR guiding principles for scientific data management and stewardship. \textit{Scientific Data}, v. 3, 2016. DOI: \url{https://doi.org/10.1038/sdata.2016.18}.
\bibitem{miller1990} MILLER, B. P.; FREDRIKSEN, L.; SO, B. An empirical study of the reliability of UNIX utilities. \textit{Communications of the ACM}, v. 33, n. 12, p. 32--44, 1990. DOI: \url{https://doi.org/10.1145/96267.96279}.
\bibitem{orchestrator2025} BECKENBAUER, L. et al. Orchestrator: active inference for multi-agent systems in long-horizon tasks. \textit{Preprint} arXiv:2509.05651, 2025. Disponível em: \url{https://arxiv.org/abs/2509.05651}. Acesso em: 10 out. 2026. (Sem DOI verificado.)
\end{thebibliography}
